% Source: https://artofproblemsolving.com/community/c3h2219004_some_identities
% Stable identity: aops-topic-2219004-q01
Since $\dfrac{i(i+1)}{2}=\dbinom{i+1}{2}$, the right side is
\[\dbinom22+\dbinom32+\cdots+\dbinom{m-1}{2}.\]
\textbf{Counting argument.} $\dbinom m3$ counts the $3$-element subsets of $\{1,2,\dots,m\}$. Sort them by their largest element $j+1$, where $2\le j\le m-1$. The other two elements are any two of $\{1,\dots,j\}$, so there are $\dbinom j2$ subsets with largest element $j+1$. Adding over $j$,
\[\dbinom m3=\sum_{j=2}^{m-1}\dbinom j2=\sum_{i=1}^{m-2}\dbinom{i+1}{2}=\sum_{i=1}^{m-2}\dfrac{i(i+1)}{2}.\qquad\blacksquare\]
(This is the hockey-stick identity: the sum of the first $m-2$ triangular numbers is a tetrahedral number. For example, $m=5$: $1+3+6=10=\dbinom53$.)
